\title{Composite Timing Metrics Do Not Measure Warning Capability:
A Video-Blind Constant Matches the Eighth of Thirteen Teams, and the
Corrected Protocol Needs Correcting}

\author{Sudaroli~Dhananjeyan,
        Surya~G,
        and~Gomathi~D%
\thanks{S. Dhananjeyan is an Independent Researcher, Bengaluru, India
(e-mail: oli.sudar@gmail.com).}%
\thanks{S. G and G. D are with T. John Institute of Technology, Bengaluru,
India.}%
\thanks{Code, all submission files and the returned scores are available at
\url{https://github.com/sudaroli1/precrash}. The measurement data are archived
at \href{https://doi.org/10.17632/yjvr5c7v3n.1}{doi:10.17632/yjvr5c7v3n.1}.}}

\markboth{IEEE Transactions on Intelligent Transportation Systems}%
{Dhananjeyan \MakeLowercase{\textit{et al.}}: Composite Timing Metrics Do Not Measure Warning Capability}

\maketitle

\begin{abstract}
Traffic-accident anticipation is increasingly scored by composite metrics that
add an earliness term --- time-to-accident at a fixed risk threshold --- to
discrimination terms such as average precision and area under the ROC curve.
The composition fails structurally: the discrimination terms are bounded on the
unit interval, the earliness terms scale with clip length, and the earliness
terms are maximised by a constant risk score above the threshold --- a
submission that, deployed, would warn continuously. We measure the consequence
twice. On a live benchmark of 1{,}417 dashcam clips and thirteen teams, whose
labels are withheld and whose metric weights are undisclosed, a constant risk
score of $0.51$ that reads no pixels scores $2.35291$: it matches the
eighth-placed team at the precision the leaderboard displays and exceeds five of
thirteen entries including our own. Twenty-five such submissions, designed so
the undisclosed terms cancel under differencing, decompose the metric without
any label: discrimination contributes $0.35105$ and the two timing terms
$2.00186$, $5.70$ times as much. The same probes establish that four curves
crossing the threshold at the same frame score $0.091$ apart, a spread the
published definition cannot produce, so the stated formula does not reproduce
its own scorer. We then validate a corrected protocol on an independent corpus
of 1{,}500 annotated clips, half containing no collision. The pathology
reproduces --- a constant takes 99.9\% of the available earliness at exactly
chance discrimination --- and the cost it hides becomes measurable: such a
warning alarms for 60 minutes in every hour of ordinary driving. The protocol
behaves as predicted but is then beaten by a control that reads only the video
file's header, because clip-level conditioning cannot see a curve that is flat
within a clip. We replace it with a frame-level statistic, coverage of the
annotated actionable window against alarm duty cycle, under which the
video-blind family scores zero and a sighted model does not. Our own
ninth-placed entry, which scores below the constant, is the case study.
\end{abstract}

\begin{IEEEkeywords}
Accident anticipation, evaluation methodology, benchmark design, video-blind
baseline, time-to-accident, collision warning systems.
\end{IEEEkeywords}
